\documentclass[12pt,a4paper]{book}
\usepackage{fontspec}
\setmainfont{Times New Roman}
\usepackage[margin=1in]{geometry}
\usepackage[hidelinks]{hyperref}
\usepackage{pdfpages}
\usepackage{setspace}

\begin{document}
\frontmatter
\begin{titlepage}
\thispagestyle{empty}
\begin{center}
{\Large Ontology of Continua}\\[0.6cm]
{\huge \textbf{Core 1.3}}\\[0.3cm]
{\Large Deutsche Navigationsausgabe der Monographie}\\[1.2cm]

{\Large Alexander Yashin}\\[0.2cm]
{\normalsize Independent Researcher}\\[0.2cm]
{\normalsize ORCID 0009-0008-6166-0914}\\[1.0cm]

{\normalsize Version 1.3}\\[0.2cm]
{\normalsize April 13, 2026}\\[1.2cm]

\begin{minipage}{0.88\textwidth}
\small
\textbf{Status dieser Ausgabe.} Dies ist die deutsche Navigationsausgabe der Leitmonographie
OC Core 1.3. Sie liefert einen lokalisierten Einstieg, erklart die Architektur des
Beweiskorpers, markiert die Claim-Grenzen, liefert eine lokalisierte theorem roadmap und bindet anschliessend den vollstandigen
englischen Masterband als kanonischen theorem-bearing Hauptkorper ein.

\medskip

\textbf{Warum diese Form.} Im aktuellen Release-Schritt bleibt der kanonische,
direkt aus dem offentlichen LaTeX-Quellkorpus gebaute Beweiskorper der englische Master.
Die deutsche Ausgabe soll deshalb vor allem eines leisten: dem deutschsprachigen Leser einen
klaren narrativen Einstieg geben, ohne die wissenschaftliche Genauigkeit der eigentlichen
Beweiskette zu verwassern.
\end{minipage}

\vfill

{\small Open foundation route: github.com/alexanderyashin/ontology-of-continua-core-main}\\
{\small Archival predecessor witness: Zenodo DOI 10.5281/zenodo.17903912}
\end{center}
\end{titlepage}

\chapter*{Zusammenfassung}
\addcontentsline{toc}{chapter}{Zusammenfassung}
\onehalfspacing
Core 1.3 ist die source-first Rekonstruktion und publikationstaugliche Hartung des
monographischen Ontology-of-Continua-Korpus. Die wissenschaftliche Aufgabe dieser Fassung
ist doppelt: Einerseits soll der volle formale Shell von Core 1.2 erhalten bleiben,
andererseits werden source witnessing, theorem-fate separation, chronology-aware repair notes,
Claim-Grenzen und outward-ready scholarly packaging explizit gemacht.

Diese Ausgabe ersetzt den englischen theorem-bearing Masterband nicht und tarnt ihn auch
nicht als kurzes Packet. Ihre Aufgabe besteht darin, dem deutschsprachigen Leser einen
klaren Einstieg zu geben: Was wird hier genau bewiesen, wie ist die Leiter aus Satzen und
Lemmata aufgebaut, warum wird die Hierarchie als `K0…K12` fixiert, wo endet der positive
bewiesene Kern und wo beginnen begrenzte Deutungen, domanenubergreifende Projektionen und
weitere Prufprogramme.

\tableofcontents

\mainmatter

\chapter{Wie diese Ausgabe zu lesen ist}
Die Monographie OC Core 1.3 ist gleichzeitig fur drei Lesergruppen gebaut.

\begin{enumerate}
\item \textbf{Editoren und Gutachter} sollen sofort erkennen konnen, was der theorem-bearing
Hauptkorper ist, wo die Beweise liegen, welche Grenzen explizit gesetzt werden und welche
Theoremfamilien ausserhalb des positiven Kerns bleiben.
\item \textbf{Formale Leser} sollen durchgehend Zugriff auf Axiome, Definitionen, Satze,
Lemmata, Formeln und Anhange haben, ohne die Struktur aus internen Ledger-Dateien
rekonstruieren zu mussen.
\item \textbf{Wissenschaftlich interessierte Nicht-Spezialisten} sollen einen klaren
narrativen Einstieg bekommen und verstehen, warum jede grosse Kapitelgruppe uberhaupt
vorhanden ist.
\end{enumerate}

Diese Ausgabe tut deshalb zwei Dinge. Zuerst liefert sie eine lokalisierte Lesekarte,
eine theorem roadmap und worked examples.
Danach bindet sie den vollstandigen englischen Masterband ein, der den kanonischen
Beweiskorper der aktuellen Releasefassung darstellt.

\chapter{Was in Core 1.3 tatsachlich bewiesen wird}
Im positiven Kern des Releases steht weiterhin die begrenzte Grammatik von Kollaps,
Residuum und Wiedergeburt. Das bedeutet:

\begin{itemize}
\item Wenn unter dem aktuellen Einbettungs- und Schwellenregime keine nichtleere zulassige
Realisierung mehr verbleibt, geht das System in einen lawful collapse state uber.
\item Eine Spur nach dem Kollaps kann als Residuum bestehen bleiben.
\item Eine spatere zulassige Fortsetzung gilt nicht automatisch als dieselbe Identitat; sie
wird als Wiedergeburt klassifiziert und unterliegt eigener Identitatsdisziplin.
\end{itemize}

Core 1.3 verstarkt daher vor allem die Disziplin, nicht die Rhetorik. Der surviving
source-native kernel wird vom ausgeschlossenen Theoremmaterial getrennt, theorem fate wird
explizit fixiert, source-audit layers werden sichtbar gemacht, und Summary-Schichten durfen
keinen internen Optimismus mehr als offentlichen Beweis ausgeben.

\chapter{Warum die Hierarchie als \texorpdfstring{$K0 \dots K12$}{K0...K12} fixiert wird}
Core 1.3 beseitigt die fruhere Unklarheit, in der manche Texte von `K0–K10` sprachen,
wahrend die tragenden source surfaces und abhangigkeitstragenden Module bereits mit einer
langeren train-Struktur arbeiteten. Im finalen Release-Doctrine ist die Hierarchie deshalb
als `K0…K12` fixiert.

\begin{itemize}
\item `K0…K10` tragen weiterhin die basale structural ladder und den Mechanismus wachsener
Ausdruckskraft.
\item `K11` wird fur rekursive meta-theoretische Achsen benotigt, auf denen nicht nur
Objekte und Ubergange, sondern auch die explanatory regimes selbst rekonstruiert werden.
\item `K12` wird fur globale Koharenzachsen benotigt, bei denen nicht die lokale
Ausdrucksstarke eines Einzelniveaus, sondern die Abstimmung der gesamten train-Struktur als
Einheitsobjekt im Mittelpunkt steht.
\end{itemize}

Die zusatzlichen Ebenen sind also nicht dekorativ. Sie sind notwendig, damit die Hierarchie
nicht genau dort abbricht, wo die Theorie bereits uber meta-level closure und whole-structure
coherence spricht.

\chapter{Karte der vollen Monographie}
Der vollstandige englische Masterband, der im Anschluss eingebunden ist, ist als
opus-magnus source zu lesen. Seine Grossstruktur ist:

\begin{enumerate}
\item Das Front Matter setzt publication role, reader promise, abstract und das Ziel der
Arbeit.
\item Die Hauptkapitel entfalten das formale Modell, embedding spaces, axioms, operator
logic und den zentralen Theorembestand.
\item Weitere Kapitel entfalten collapse / death / residue / rebirth semantics,
cross-level constraints und die domanenbezogene Interpretationsgrenze.
\item Reader guide, source audit und foundational consistency chapters liefern genau das,
was Core 1.3 zusatzlich braucht: eine klare didactic spine und explizite publication
discipline.
\item Die Anhange halten provenance, source audit, extraction boundary between master
monograph and journal core sowie zusatzliche support materials fest.
\end{enumerate}

\chapter{Claim-Grenze und was bewusst nicht behauptet wird}
Der wissenschaftliche Wert dieser Monographie hangt nicht an rhetorischem Maximalismus,
sondern an strenger Trennung der Aussageklassen.

\begin{itemize}
\item \textbf{Axiomatische Aussagen} setzen die basalen formalen Voraussetzungen.
\item \textbf{Abgeleitete Aussagen} gehoren nur dann zum theorem-bearing Korpus, wenn
dependency chain und proof route explizit nachvollziehbar bleiben.
\item \textbf{Strukturelle Deutungen} sind als Erklarungsschichten zulassig, aber nicht
automatisch neue Theoreme.
\item \textbf{Domänenubergreifende Projektionen} dienen als Karte weiterer Arbeit, nicht als
offentlicher Beweis einer bereits vollzogenen Gesamtschliessung.
\end{itemize}

Deshalb vermeidet Core 1.3 bewusst jede Formulierung, die uber den tatsachlichen proof
surface hinausgeht. Die Ambition des Frameworks bleibt sichtbar, doch das public wording
bleibt sober, exact and conventionally scientific.

\chapter{Wie sich der Journal Core dazu verhalt}
Der Journal Core ersetzt die Monographie nicht. Er ist eine traceable extraction from the
master. Seine Funktion besteht darin, Editoren und Gutachtern einen kurzeren editorial route
anzubieten, ohne die kanonische Bindung an den grossen Beweiskorper zu verlieren.

Die sinnvolle Lesereihenfolge lautet daher:

\begin{enumerate}
\item Wer den vollen Beweis- und Didaktik-Kontext braucht, beginnt mit der
Master-Monographie.
\item Wer zuerst den kompakten Editorial-Route fur ein Journal braucht, kann mit dem
Journal Core beginnen, muss ihn aber als Extract und nicht als eigenstandige Theorie lesen.
\item Alle companion packets sind als support layers um diese beiden science artifacts zu
verstehen, nicht als Ersatz der Manuskripte.
\end{enumerate}

\chapter{Der kanonische englische Korpus}
Im Folgenden ist der vollstandige englische Masterband eingebunden, built directly from the
public LaTeX source corpus. Genau dieser Band bildet den canonical theorem-bearing body der
Releasefassung Core 1.3.

\cleardoublepage
\includepdf[pages=-]{build_oc_core_1_3_master/oc_core_1_3_master_monograph.pdf}

\end{document}
